% 第1章：项目概述。
%
% 本章沿用初赛文档的叙述方式，在原有技术背景上更新多核、双架构和用户态工作负载。

\chapter{项目概述}

\wateros{}是用Rust编写的\code{no\_std}操作系统内核，面向操作系统竞赛在QEMU上完成
bring-up与测例验收。内核已经在RISC-V64 OpenSBI与LoongArch64 \code{virt}两条主线上
跑通启动、trap分发、分页、任务调度与用户ELF执行；文件系统、VFS、系统调用、IPC和
网络也接在同一套用户态路径上。在原有单核主线的基础上，当前实现进一步接通了八核
启动、per-CPU运行状态、核间中断、跨核任务唤醒和TLB shootdown。

两条平台路径的入口参数、trap frame、页表格式、定时器与VirtIO总线不同，差异收束在
\code{wateros-platform}、\code{wateros-driver}和\code{wateros-mm}各自的平台或架构实现中。
进入task、FS/VFS、IPC、network、cred、klog和syscall之后，上层逻辑共用同一套接口，
用户程序在两条路径上看到的fd、进程、地址空间与系统调用编号保持一致。多核支持也沿用
这条边界：CPU启动、IPI和本地TLB操作由架构层完成，任务状态、地址空间生命周期和缓存
一致性仍由对应组件负责。

立项时除了竞赛验收，我们也希望内核能支撑多人并行开发和教学演示。单体仓库里改一处
驱动往往要牵一串头文件和全局状态；我们更想要的是：组员互不了解对方模块细节，仍能在
抽象层上对接。熟悉某一域的人把\code{api-v0}里的trait和类型定下来，再用本层原语往上
拼接口；只打算贡献某一个\code{impl-*}的人可以在crate内部实现细节，不必追问这些原语
最后被谁、以何种方式调用。组装操作系统时因此主要面对稳定模块名和契约，而不是同时
盯着所有底层实现。

Rust和这种拆法比较合拍。\code{api-v0}用trait列出必须实现的行为，\code{pub}边界控制
后来接手的人能看到的面；所有权和借用检查把一部分内存与并发问题拦在编译期，代码评审
时可以把精力放在逻辑是否正确、feature组合是否合理。每个\code{impl-*}在
\code{Cargo.toml}里注册成feature之后，根crate切换feature即可更换实现。驱动允许多个
实现按设备类型组合；堆分配器、平台arch实现和调度器策略这类互斥选项，则在聚合层检查
配置冲突，避免一组无法工作的实现进入内核镜像。

目前的开发重点已经从“让单个测例进入内核”转向“让完整用户态工作流稳定结束”。除原有
basic、BusyBox、glibc/musl和LTP定向用例外，内核还需要承受CAgent中的并发网络请求，
以及BuildStorm中Cargo、rustc、多进程、多线程、动态链接、内存映射和大量文件读写的
组合压力。这些负载没有引入另一套特化接口，而是直接复用现有的task、MM、VFS、网络和
系统调用主路径，用来暴露单项测例不容易触发的并发与资源回收问题。

\section{设计目标}

\begin{itemize}
  \item 在QEMU中完成可启动、可抢占调度、可装载并执行用户ELF的内核主线，并在
  RISC-V64与LoongArch64上支持八核启动和多核任务运行。
  \item 系统调用编号与参数布局对齐Linux generic 64-bit ABI，使根文件系统中的
  glibc/musl程序能够沿标准系统调用路径进入内核。
  \item 在两种架构上分别接通VirtIO块设备与网卡、ext4读写根卷和VFS fd会话，支撑
  BusyBox脚本、网络服务以及Cargo/rustc工具链工作负载。
  \item 保持平台相关代码与MM、task、VFS、IPC、network、syscall等上层模块分离，
  新增平台或替换具体实现时不改动上层聚合接口。
  \item 明确多核执行中的任务唯一运行、锁边界、跨核唤醒、TLB失效和对象回收条件，
  使并发正确性不依赖偶然的运行时序。
\end{itemize}

\section{当前实现摘要}

\begin{longtable}{L{3.2cm}L{10.3cm}}
  \caption{WaterOS当前实现摘要}\label{tab:current-summary}\\
  \toprule[1.5pt]
  子系统 & 当前实现状态 \\
  \midrule[1pt]
  \endfirsthead
  \toprule[1.5pt]
  子系统 & 当前实现状态 \\
  \midrule[1pt]
  \endhead
  启动与平台 &
  RISC-V64从OpenSBI接收hart id与DTB物理地址，通过SBI HSM启动从核；LoongArch64使用
  QEMU \code{virt}平台的启动、异常和核间中断路径。BSP负责全局服务初始化，AP使用独立
  启动栈完成本核trap、timer、CPU-local和idle task初始化，全部就绪后进入online集合。\\
  内存管理 &
  RISC-V64使用Sv39页表，LoongArch64使用独立的三级4 KiB页表实现。当前已接通内核页表、
  用户地址空间、ELF装载、COW fork、惰性mmap、用户缺页和user access拷贝；页表修改会
  根据地址空间的活动CPU集合执行本地失效或跨核TLB shootdown。\\
  任务与调度 &
  task组件提供内核与用户任务、fork/clone/exec、线程组、wait族和zombie回收。调度状态
  按CPU组织，任务迁移、阻塞唤醒、timer tick和reschedule IPI共同维持多核运行；任务
  退出后由统一的生命周期状态决定何时从调度结构移除并释放资源。\\
  IPC与信号 &
  已实现waitqueue、pipe、futex、SysV共享内存和信号主路径。futex等待与唤醒和任务状态
  转换相连，信号覆盖disposition、pending、mask、signal frame与基础的syscall restart，
  并处理多线程退出、组退出和父进程等待之间的关系。\\
  驱动 &
  RISC-V64经设备树扫描VirtIO-MMIO块设备与网卡，LoongArch64经PCI枚举VirtIO设备。
  块设备注册到统一设备表供文件系统使用；网络设备接入smoltcp协议栈；字符设备提供
  控制台及用户态所需的设备节点。\\
  文件系统与VFS &
  ext4读写根卷通过FS接口接入VFS，devfs、procfs和ramfs挂在同一命名空间内。VFS提供
  mount table、页缓存、路径解析和per-task fd会话，支持cwd、dup、fork继承与CLOEXEC；
  写路径需要同时维护页缓存、inode大小、目录项和底层ext4元数据的一致性。\\
  网络 &
  socket系统调用通过统一fd对象连接到smoltcp及UNIX socket实现，覆盖CAgent所需的
  TCP/UDP并发连接、poll类等待和常用socket选项。网络poller、设备队列和用户线程之间
  使用明确的唤醒关系，避免只能依赖定时轮询推进连接状态。\\
  系统调用与ABI &
  Linux generic64系统调用层覆盖文件IO、进程线程、内存、时间、目录、凭证、pipe、
  poll/epoll、futex、SysV IPC、INET/UNIX socket等主路径。返回值统一使用Linux错误码，
  用户指针在进入具体子系统前完成范围、权限和跨页检查。\\
  验证与复现 &
  原有basic、BusyBox、glibc/musl和LTP定向入口继续用于回归；CAgent检查网络、计时和
  系统状态接口，BuildStorm检查Cargo/rustc并行编译链。两种架构分别保留构建、启动、
  客体成功标记和文件系统检查记录，不以宿主脚本退出码代替用户态结果。\\
  \bottomrule[1.5pt]
\end{longtable}

\section{当前局限}

本节记录当前实现与完整Linux语义之间仍存在的边界。这些限制不影响已经跑通的主路径，
但在增加用户程序或扩大压力测试时仍需继续验证。

\begin{itemize}
  \item \strong{Linux语义仍是子集。}复杂job control、少用系统调用、完整capability、
  部分文件系统flag和极端信号组合仍有未覆盖边界，不能把测例通过等同于完整Linux兼容。
  \item \strong{双架构需要分别验证。}两种架构共用上层机制，但固件入口、页表、TLB、
  中断和VirtIO总线不同；一种架构上的运行结果不能代替另一种架构。
  \item \strong{多核放大了生命周期问题。}任务、地址空间、fd、socket和缓存对象可能
  同时被多个CPU访问，退出、取消映射和回收路径仍需通过更长时间的并发负载继续检查。
  \item \strong{文件系统异常路径覆盖有限。}当前验证以正常运行、同步写和正常关闭后的
  完整性检查为主，掉电点注入、journal恢复和复杂多挂载场景仍需要单独测试。
  \item \strong{运行环境需要固定。}完整工作负载依赖匹配版本的QEMU、固件、根文件系统
  和用户态工具链。提交前需要在同一版本上重新运行，避免混用不同阶段的日志和结果。
\end{itemize}
